% Options for packages loaded elsewhere
\PassOptionsToPackage{unicode}{hyperref}
\PassOptionsToPackage{hyphens}{url}
\PassOptionsToPackage{dvipsnames,svgnames,x11names}{xcolor}
%
\documentclass[
  11pt,
  a4paper,
]{ltjsarticle}
\usepackage{amsmath,amssymb}
\usepackage{iftex}
\ifPDFTeX
  \usepackage[T1]{fontenc}
  \usepackage[utf8]{inputenc}
  \usepackage{textcomp} % provide euro and other symbols
\else % if luatex or xetex
  \usepackage{unicode-math} % this also loads fontspec
  \defaultfontfeatures{Scale=MatchLowercase}
  \defaultfontfeatures[\rmfamily]{Ligatures=TeX,Scale=1}
\fi
\usepackage{lmodern}
\ifPDFTeX\else
  % xetex/luatex font selection
\fi
% Use upquote if available, for straight quotes in verbatim environments
\IfFileExists{upquote.sty}{\usepackage{upquote}}{}
\IfFileExists{microtype.sty}{% use microtype if available
  \usepackage[]{microtype}
  \UseMicrotypeSet[protrusion]{basicmath} % disable protrusion for tt fonts
}{}
\makeatletter
\@ifundefined{KOMAClassName}{% if non-KOMA class
  \IfFileExists{parskip.sty}{%
    \usepackage{parskip}
  }{% else
    \setlength{\parindent}{0pt}
    \setlength{\parskip}{6pt plus 2pt minus 1pt}}
}{% if KOMA class
  \KOMAoptions{parskip=half}}
\makeatother
\usepackage{xcolor}
\usepackage[margin=24mm]{geometry}
\setlength{\emergencystretch}{3em} % prevent overfull lines
\providecommand{\tightlist}{%
  \setlength{\itemsep}{0pt}\setlength{\parskip}{0pt}}
\setcounter{secnumdepth}{5}
\ifLuaTeX
  \usepackage{selnolig}  % disable illegal ligatures
\fi
\IfFileExists{bookmark.sty}{\usepackage{bookmark}}{\usepackage{hyperref}}
\IfFileExists{xurl.sty}{\usepackage{xurl}}{} % add URL line breaks if available
\urlstyle{same}
\hypersetup{
  colorlinks=true,
  linkcolor={blue},
  filecolor={Maroon},
  citecolor={Blue},
  urlcolor={blue},
  pdfcreator={LaTeX via pandoc}}

\author{}
\date{}

\begin{document}

\hypertarget{chapter-11-ux6570ux5024ux7ddaux5f62ux4ee3ux6570}{%
\section{Chapter 11 ---
数値線形代数}\label{chapter-11-ux6570ux5024ux7ddaux5f62ux4ee3ux6570}}

数学上正しい式が、コンピュータ上で良い計算法とは限りません。この章ではその差を扱います。

\hypertarget{ux4e8cux3064ux306eux8aa4ux5deeux6e90ux3092ux5206ux3051ux308b}{%
\subsection{二つの誤差源を分ける}\label{ux4e8cux3064ux306eux8aa4ux5deeux6e90ux3092ux5206ux3051ux308b}}

\hypertarget{ux554fux984cux306e-conditioning}{%
\subsubsection{問題の
conditioning}\label{ux554fux984cux306e-conditioning}}

入力が少し変わったとき、真の答えがどれだけ変わるか。

\hypertarget{ux30a2ux30ebux30b4ux30eaux30baux30e0ux306e-stability}{%
\subsubsection{アルゴリズムの
stability}\label{ux30a2ux30ebux30b4ux30eaux30baux30e0ux306e-stability}}

有限精度計算によって、アルゴリズム自身がどれだけ余計な誤差を生むか。

この二つは全く別です。

\hypertarget{svd-ux3067-conditioning-ux3092ux898bux308b}{%
\subsection{SVD で conditioning
を見る}\label{svd-ux3067-conditioning-ux3092ux898bux308b}}

\[
\kappa_2(A)=\frac{\sigma_{max}}{\sigma_{min}}
\]

は、線形変換が最もよく見える方向と最も見えにくい方向の比です。

最小特異値が小さいほど、逆問題ではその方向のノイズが大きく増幅されます。

\hypertarget{ux306aux305c-inva-b-ux3092ux907fux3051ux308bux306eux304b}{%
\subsection{\texorpdfstring{なぜ \texttt{inv(A)\ @\ b}
を避けるのか}{なぜ inv(A) @ b を避けるのか}}\label{ux306aux305c-inva-b-ux3092ux907fux3051ux308bux306eux304b}}

理論では

\[
x=A^{-1}b
\]

でも、数値計算では明示的な逆行列を作らず LU/QR などで \texttt{Ax=b}
を直接解く方が通常効率的かつ安定です。

\hypertarget{ux5927ux898fux6a21ux554fux984c}{%
\subsection{大規模問題}\label{ux5927ux898fux6a21ux554fux984c}}

行列が巨大になると完全な SVD
や固有値分解は不要・不可能なことがあります。そこで上位成分だけを求める
Krylov 法や randomized SVD を使います。

\hypertarget{ux5230ux9054ux57faux6e96}{%
\subsection{到達基準}\label{ux5230ux9054ux57faux6e96}}

同じ数学的解を返す複数の式を見たとき、「どれを計算機で使うべきか」を
conditioning・計算量・stability の観点から考えられることを目標にします。

\end{document}
